\chapter{Competitor Analysis}

\section{The Competitive Field}

AMR players fall into groups:
\begin{itemize}
    \item Warehousing/manufacturing: MiR, Geek+, Locus, Hai, Exotec, Fetch/Zebra; Robotnik for R\&D \cite{standardbots_amr,fdatabot_amr2026,robotnik_kairos}.
    \item Inspection: Boston Dynamics (Spot), Capra, DEEP \cite{standardbots_amr,generationrobots_amr}.
    \item Security: Knightscope, Cobalt, SMP; Tunisian Enova Robotics \cite{mordor_security,coherent_security}.
    \item Healthcare/hospitality: Aethon TUG, Savioke (purpose-built).
    \item Retail: Bossa Nova \cite{mir_healthcare,nordmodules_retail}.
\end{itemize}

Table~\ref{tab:competitors} provides a structured comparison.

\section{Cross-Competitor Patterns} 
 
Four patterns emerge:

\textbf{1. Modularity already exists at the top.} MiR's base accepts many attachments. The differentiator must be price, local support, and ease of deployment on unmodified sites -- not modularity per se.

\textbf{2. Shelf access splits; the harder part is crowded.} Vertical reach specialists (Exotec, Hai, etc.) require engineered racking, not retrofits. Retrofit is risky and competes against well-funded firms -- defer it.

\textbf{3. Unit economics are unforgiving; capex-light models win.} Even leaders like Locus have struggled (Chapter 11). RaaS is the way to go for a Tunisian entrant.

\textbf{4. All competitors target high-purchasing-power markets.} None are adapted to Tunisia's cost-to-labor ratio. They lack local manufacturing, support, and French/Arabic presence -- this is the gap.

% Competitor comparison table (referenced as \ref{tab:competitors}).
% Fitted to text width: footnotesize, ragged columns so long words wrap
% instead of crossing the cell border. Content unchanged.
{\footnotesize
\setlength{\tabcolsep}{3pt}
\renewcommand{\arraystretch}{1.3}
\begin{longtable}{|>{\raggedright\arraybackslash}m{1.7cm}|
                  >{\raggedright\arraybackslash}m{1.6cm}|
                  >{\raggedright\arraybackslash}m{2.0cm}|
                  >{\raggedright\arraybackslash}m{2.5cm}|
                  >{\raggedright\arraybackslash}m{1.7cm}|
                  >{\raggedright\arraybackslash}m{1.8cm}|
                  >{\raggedright\arraybackslash}m{2.2cm}|
                  >{\centering\arraybackslash}m{1.4cm}|}
\caption{Competitor comparison -- global and regional AMR platforms.}\label{tab:competitors}\\
\hline
\rowcolor{gray!15}
\textbf{Company} & \textbf{Product} & \textbf{Sector} & \textbf{Key features} & \textbf{Pricing} & \textbf{Strengths} & \textbf{Weaknesses} & \textbf{Source} \\
\hline
\endfirsthead
\hline
\rowcolor{gray!15}
\textbf{Company} & \textbf{Product} & \textbf{Sector} & \textbf{Key features} & \textbf{Pricing} & \textbf{Strengths} & \textbf{Weaknesses} & \textbf{Source} \\
\hline
\endhead
\hline
\endfoot
MiR (Teradyne) & MiR250 / 600 / 1350 & Warehousing, Mfg. & Interchangeable top modules; arm option (MC600) & USD 30--45k & Mature, large install base & Proprietary software lock-in & \cite{mir250} \\
\hline
Robotnik & RB-KAIROS+ & Mfg., R\&D & ROS2 base, plug-and-play arm & From EUR 50k & Open, modular & High entry price & \cite{robotnik_kairos} \\
\hline
Geek+ & P-series, RoboShuttle & Warehousing & Goods-to-person, sortation, shuttle & Project-based & Huge fleet scale & System-level, not one robot & \cite{fdatabot_amr2026} \\
\hline
Locus Robotics & LocusBots & Warehousing, 3PL & Collaborative picking, RaaS & RaaS subscription & Capex-light model & Single workflow; financial stress & \cite{standardbots_amr,builtin_raas} \\
\hline
Hai Robotics & HaiPick / Climb & Warehousing (totes) & Tote/case ACR, vertical reach & Project-based & Vertical density & Needs engineered racking & \cite{standardbots_amr} \\
\hline
Exotec & Skypod & Warehousing (high-density) & Shelf-climbing goods-to-person & Project, multi-million & High throughput & System-level only, costly & \cite{fdatabot_amr2026} \\
\hline
Boston Dynamics & Spot & Inspection & Legged; camera/sensor mounts & USD 75k+ & Best-in-class mobility & Very high cost and training & \cite{standardbots_amr} \\
\hline
Capra Robotics & Capra Hircus & Outdoor inspection/patrol & Fully modular, rugged & Quote-based & Outdoor-rugged & Narrow niche & \cite{generationrobots_amr} \\
\hline
Fetch / Zebra & Freight AMRs & Warehouse transport & Transport + data & Project / fleet & Zebra channel reach & Transport-focused only & \cite{standardbots_amr} \\
\hline
Aethon & TUG & Healthcare, hospitality & Secure cart delivery & Quote-based subscription & Proven in hospitals & Single-purpose, not reconfigurable & \cite{researchgate_hospital} \\
\hline
Cobalt Robotics & Indoor patrol & Security & Patrol + teleoperation & USD 3.5--6.5k/mo & Indoor security & Security-only & \cite{coherent_security} \\
\hline
SMP Robotics & Outdoor patrol & Security & Additional sensors & USD 120--320k & Large-area patrol & Extremely high cost & \cite{mordor_security} \\
\hline
Bossa Nova & Shelf scanner & Retail & Scanner only & Undisclosed & Shelf analytics & Scanning only; Walmart ended use & \cite{nordmodules_retail,dudarenok_retail_china} \\
\hline
Enova Robotics (TN) & PGuard & Security, inspection & Custom design + AI; locally built & Undisclosed & Local capability; exported & UGV/security-focused & TBD (local) \\
\hline
\end{longtable}
}

\newpage